\scanpage{37}
\mspage{11}
En résumé et sauf erreur

\uline{Prop.} Soit $X$ une sphère conforme \deleted{\unsure{orientée}}, $H\subset\operatorname{Aut}_{\mathrm{conf}}(X)\overset{\mathrm{def}}=G$ un sous-groupe \uline{compact} \deleted{de}, Alors

a) Il existe une métrique riemannienne sphérique normalisée sur $X$ comp. avec sa structure conforme

b) Celle-ci est unique ssi \deleted{il n'existe pas de}\inserted{$X^H=\varnothing$}

\deleted{\unsure{d}} c) Supposons \deleted{\unsure{$X^H$}} $X^H\ne\varnothing$, $H\ne\{1\}$. Alors On a \deleted{\unsure{$X$}}\inserted{ou bien} $\operatorname{card}X^H=2$, ou bien $H=\{1,\sigma\}$ où $\sigma$ est une inversion.

1) Si $\operatorname{card}X^H=2$, \deleted{Soit} $H$\inserted{$^+=H\cap G^+$. On a $X^H=X^{H^+}$, et $H^+$} est contenu dans un unique tore compact $T$, et \deleted{\unsure{les métriques}}\inserted{les métriques} sphériques riemanniennes normalisées comp., invariantes par \deleted{\unsure{$T$}} $H$, sont en corr. biunivoque avec les orbites de\inserted{$T$ dans} $(X\setminus X^H$\deleted{$)/T$} [donc tte orbite de \deleted{\unsure{$T$}} $H^+$ dans $X\setminus X^H$ définit une telle métrique]

\deleted{\unsure{b}} 2) Si $H=\{1,\sigma\}$, $\sigma$ une inversion, les métriques invariantes sous $H$ sont en corr. biunivoque avec les éléments de
\[
(X\setminus X^H)/H=(X\setminus X^\sigma)/\sigma
\]

\margin{\uline{à \unsure{rectifier|vérifier}} (\uline{NB} Il faudrait arriver à traduire en termes d'alg. linéaire, l'action d'un autom. $\unsure{u|h}\in G$ sur une métrique sphérique\unsure{.})}
% archive foliation: 39
